%% Section 3: Taxonomy of Repository Representation Paradigms
\section{Taxonomy of Repository Representation Paradigms}
\label{sec:taxonomy}

We identify seven \emph{Paradigm of Automated Repository Representation} (PAR) categories,
differentiated along ten structural dimensions: (D1) granularity of retrieval unit,
(D2) index structure, (D3) retrieval mechanism, (D4) context assembly strategy,
(D5) cross-file dependency modelling, (D6) dynamic context update, (D7) multi-modal
artefact support, (D8) downstream task scope, (D9) evaluation benchmark, and
(D10) training dependency (Zero-shot~L0, RAG-only~L1, Lightweight FT~L2, Full RL~L3).

\subsection{PAR-1: Dense Embedding Retrieval}
\label{subsec:par1}

PAR-1 systems encode repository artefacts as fixed-size dense vectors and retrieve
relevant snippets via approximate nearest-neighbour search.  Representative systems
include CodeBERT~\cite{feng2020codebert}, GraphCodeBERT~\cite{guo2021graphcodebert},
and UniXcoder~\cite{guo2022unixcoder}.

\subsection{PAR-2: Retrieval-Augmented Generation (Cross-File)}
\label{subsec:par2}

PAR-2 systems perform query-driven retrieval from a file-level index and prepend
retrieved chunks to the LLM prompt.  RepoCoder~\cite{zhang2023repocoder},
Repoformer~\cite{wu2024repoformer}, ReACC~\cite{lu2022reacc}, and
RepoGenix~\cite{liang2024repogenix} exemplify this paradigm.

\subsection{PAR-3: Graph-Based Structural Representation}
\label{subsec:par3}

PAR-3 systems build explicit structural graphs: Abstract Syntax Trees (ASTs), Control
Flow Graphs (CFGs), Call Graphs, Data Flow Graphs (DFGs), or Knowledge Graphs.
GraphCoder~\cite{liu2024graphcoder}, RepoGraph~\cite{ouyang2025repograph},
RepoHyper~\cite{phan2025repohyper}, and CodexGraph~\cite{liu2025codexgraph} belong
to this paradigm.  GraphCodeBERT~\cite{guo2021graphcodebert} integrates a DFG
during pre-training and straddles PAR-1/PAR-3.

\subsection{PAR-4: Execution-Trace and Dynamic Context}
\label{subsec:par4}

PAR-4 systems derive context from runtime artefacts (stack traces, test output,
compiler errors) or passive static-analysis traces rather than constructing an
explicit offline index.  SWE-agent~\cite{yang2024sweagent} and
AutoCodeRover~\cite{zhang2024autocoderover} exemplify reactive, execution-driven
context assembly.

\subsection{PAR-5: Memory-Augmented Representation}
\label{subsec:par5}

PAR-5 systems maintain a persistent, mutable memory store that \emph{is} the context
window for a given episode.  CodeMEM~\cite{wang2026codemem} uses AST-guided adaptive
memory; SWE-Exp~\cite{chen2025sweexp} accumulates experience traces across issues.
Unlike PAR-7 (where memory is one component of a broader agent loop), in PAR-5 the
memory store is the primary representation substrate.

\subsection{PAR-6: Compression and Summarisation}
\label{subsec:par6}

PAR-6 systems reduce token footprint by summarising or compressing repository content
before passing it to the LLM.  LongCodeZip~\cite{shi2025longcodezip} learns
differentiable soft-token compression; hierarchical pruning
approaches~\cite{zhang2025hierarchicalpruning,dhulshette2025hiersum} and
SWE-Pruner~\cite{wang2026swepruner} apply hierarchical or agent-driven summarisation.

\subsection{PAR-7: Agentic Repository Navigation}
\label{subsec:par7}

PAR-7 systems employ an autonomous agent loop—tool calls, file reads, searches,
code execution—to \emph{navigate} the repository on demand.  Memory is one component
among many.  CodePlan~\cite{bairi2024codeplan}, MetaGPT~\cite{hong2023metagpt},
OpenHands~\cite{wang2025openhands}, LingmaAgent~\cite{ma2025lingmaagent}, and
Prometheus~\cite{pan2026prometheus} belong to this paradigm.
